\section{Εκτέλεση χωρίς τοπική διατήρηση κατάστασης και αποθήκευση σε εξωτερικά αποθετήρια}
\label{sec:stateless_externalized_state}

Η υλοποίηση ακολουθεί συχνά το μοτίβο εκτέλεσης χωρίς τοπική διατήρηση κατάστασης (\ENG{stateless runtime}), αλλά η διατύπωση αυτή χρειάζεται ακρίβεια. Ο \ENG{engine} δεν διατηρεί στη μνήμη του \ENG{application process} την έγκυρη κατάσταση της συνεδρίας. Η ενεργή κατάσταση κάθε \ENG{launch} αποθηκεύεται σε εξωτερικό αποθετήριο στη \ENG{Redis} μέσω του \texttt{\ENG{RuntimeStateStore}}. Έτσι η κατάσταση των χρηστών είναι ρητά κατανεμημένη σε εξειδικευμένα συστήματα αποθηκεύσης με διαφορετικό βαθμό μονιμότητας και διαφορετικό ρόλο ως προς την έγκυρη μορφή αναφοράς.

Το αντικείμενο \texttt{\ENG{LaunchRuntimeState}} συγκεντρώνει τις κρίσιμες πληροφορίες μιας ενεργής συνεδρίας: \ENG{launchId}, \ENG{attemptId}, χρήστη, μάθημα, πρότυπο, τελευταία αποδεκτή ακολουθία, τελευταία χρονική στιγμή εμφάνισης και το \texttt{\ENG{NormalizedProgress}}. Ο \texttt{\ENG{LaunchRuntimeStateFactory}} το αρχικοποιεί κατά τη δημιουργία του \ENG{launch}, ενώ ο \texttt{\ENG{RuntimeOrchestrator}} το ενημερώνει σε κάθε αποδεκτό \ENG{commit}. Το \texttt{\ENG{RedisRuntimeStateStore}} σειριοποιεί το αντικείμενο ως \ENG{JSON}, ορίζει \ENG{TTL} και επιτρέπει τόσο την ανάκτηση μιας συγκεκριμένης κατάστασης όσο και τον εντοπισμό παλαιών, μη συγχρονισμένων (\ENG{dirty}) καταστάσεων.

Η μόνιμη και έγκυρη αποθήκευση τοποθετείται στην \ENG{Postgres}. Σύμφωνα με τον σχεδιασμό της βάσης ο πίνακας \texttt{\ENG{attempts}} αποθηκεύει το συγκεντρωτικό αποτέλεσμα μιας προσπάθειας, δηλαδή \ENG{completion status}, \ENG{success status}, βαθμολογίες, συνολικό χρόνο, τελευταία θέση και τη χρονική στιγμή του τελευταίου \ENG{commit}. Παράλληλα, ο πίνακας \texttt{\ENG{attempt\_runtime\_snapshots}} διατηρεί στιγμιότυπα του \ENG{raw runtime JSON} και της κανονικοποιημένης προόδου. Η διάκριση αυτή είναι σημαντική καθώς η \ENG{Redis} διατηρεί την ενεργή, μεταβαλλόμενη κατάσταση, ενώ η \ENG{Postgres} διατηρεί την ιστορική και τελική αναπαράσταση που πρέπει να παραμένει ανακτήσιμη μετά τη λήξη της συνεδρίας.

Η μεταφορά από την προσωρινή κατάσταση στη μόνιμη αποθήκευση γίνεται από τον \texttt{\ENG{RuntimeFlushService}}. Κατά τον κανονικό τερματισμό, η υπηρεσία ανακτά το ενεργό \ENG{LaunchRuntimeState}, εφαρμόζει το \texttt{\ENG{NormalizedProgress}} στο \ENG{attempt}, αποθηκεύει στιγμιότυπο των δεδομένων και κλείνει τόσο την προσπάθεια όσο και το \ENG{launch}. Η ίδια υπηρεσία χρησιμοποιείται και από τον \texttt{\ENG{RuntimeFlushScheduler}}, ο οποίος περιοδικά εντοπίζει χρονικά παρωχημένα ενεργά \ENG{launches} και τα κλείνει ως \ENG{ABANDONED}. Έτσι, η πλατφόρμα δεν εξαρτά την τελική καταγραφή μόνο από ομαλούς τερματισμούς μέσω του \ENG{browser}.

Το Σχήμα~\ref{fig:stateless_state} αποτυπώνει αυτή την κατανομή ρόλων αναφοράς. Το Σχήμα~\ref{fig:runtime_flush_sequence} εξειδικεύει τη ροή συγχρονισμού προς μόνιμη αποθήκευση από το ενεργό \ENG{launch}. Ο Πίνακας~\ref{tab:state_authority_matrix} συμπληρώνει το σχήμα, καταγράφοντας για κάθε αποθετήριο αν λειτουργεί ως κύριο αποθετήριο αναφοράς ή ως υποστηρικτικό επίπεδο.

\begin{figure}[H]
  \centering
\begin{chapterfigurebox}
  \begin{tikzpicture}[
      font=\footnotesize,
      node distance=1.8cm and 2.8cm,
      box/.style={draw, rounded corners, fill=gray!5, align=center, text width=3.2cm, minimum height=1.08cm},
      emph/.style={draw, rounded corners, fill=gray!12, align=center, text width=3.5cm, minimum height=1.08cm, font=\footnotesize\bfseries},
      store/.style={draw, cylinder, shape border rotate=90, aspect=0.28, minimum height=1.2cm, minimum width=3.2cm, align=center, fill=gray!5},
      arrow/.style={-Latex, thick}
    ]
    \node[box] (browser) {Κατάσταση μνήμης\\στον \ENG{browser}};
    \node[emph, right=of browser] (engine) {Πυρήνας εκτέλεσης\\\ENG{engine}};
    \node[store, right=of engine] (minio) {\ENG{MinIO}\\συμπιεσμένο περιεχόμενο};

    \node[store, below=1.9cm of engine] (redis) {\ENG{Redis}\\ενεργή κατάσταση \ENG{launch}};
    \node[emph, below=1.9cm of minio] (playercache) {\ENG{Player cache}\\αποσυμπιεσμένα αρχεία};
    \node[store, below=1.8cm of redis] (postgres) {\ENG{Postgres}\\τελικές εγγραφές προσπαθειών};

    \draw[arrow] (browser) -- (engine);
    \draw[arrow] (engine) -- (redis);
    \draw[arrow] (redis) -- (postgres);
    \draw[arrow] (minio) -- (playercache);

    \node[font=\footnotesize, align=center, below=0.18cm of redis] {προσωρινή έγκυρη κατάσταση};
    \node[font=\footnotesize, align=center, below=0.18cm of postgres] {μόνιμο ιστορικό};
    \node[font=\footnotesize, align=center, below=0.18cm of playercache] {μόνο για διανομή};
  \end{tikzpicture}
\end{chapterfigurebox}
  \caption{Κατανομή της κατάστασης ανάμεσα σε \ENG{browser}, \ENG{player cache}, \ENG{Redis}, \ENG{Postgres} και \ENG{MinIO}, με σαφή διάκριση του έγκυρου σημείου αναφοράς.}
  \label{fig:stateless_state}
\end{figure}

\begin{figure}[H]
  \centering
\begin{chapterfigurebox}
  \begin{tikzpicture}[
      font=\small,
      box/.style={draw, rounded corners, fill=gray!5, align=center, text width=7.2cm, minimum height=1.0cm},
      emph/.style={draw, rounded corners, fill=gray!12, align=center, text width=7.4cm, minimum height=1.0cm, font=\small\bfseries},
      arrow/.style={-Latex, thick}
    ]
    \node[emph] (start) at (0,0) {Ενεργό \ENG{launch} στο \ENG{Redis}};
    \node[box] (trigger) at (0,-1.7) {Ενεργοποίηση: κανονικός τερματισμός ή λήξη χρονικού ορίου αδράνειας};
    \node[emph] (flush) at (0,-3.4) {\texttt{\ENG{RuntimeFlushService}}};
    \node[box] (apply) at (0,-5.1) {Εφαρμογή κανονικοποιημένης προόδου στο \texttt{\ENG{attempt}}};
    \node[box] (snapshot) at (0,-6.8) {Αποθήκευση \texttt{\ENG{attempt\_runtime\_snapshot}} με αιτιολόγηση};
    \node[box] (close) at (0,-8.5) {Κλείσιμο \texttt{\ENG{attempt}} και \texttt{\ENG{launch}}, διαγραφή κατάστασης από \ENG{Redis}};

    \draw[arrow] (start) -- (trigger);
    \draw[arrow] (trigger) -- (flush);
    \draw[arrow] (flush) -- (apply);
    \draw[arrow] (apply) -- (snapshot);
    \draw[arrow] (snapshot) -- (close);
  \end{tikzpicture}
\end{chapterfigurebox}
  \caption{Ακολουθία συγχρονισμού προς μόνιμη αποθήκευση από ενεργή κατάσταση εκτέλεσης σε τελική κατάσταση \ENG{attempt} και ιστορικό στιγμιότυπο.}
  \label{fig:runtime_flush_sequence}
\end{figure}


\begin{table}[H]
\centering
\caption{Ρόλοι αναφοράς και λειτουργίες των αποθετηρίων κατάστασης στην υλοποίηση}
\label{tab:state_authority_matrix}
\begingroup
\footnotesize
\setlength{\tabcolsep}{3pt}
\renewcommand{\arraystretch}{1.20}
\begin{adjustbox}{center,max width=\textwidth}
\begin{tabular}{|p{2.8cm}|p{2.8cm}|p{5.2cm}|p{4.0cm}|}
\hline
Αποθετήριο / τεχνούργημα & Ρόλος ως προς την έγκυρη μορφή αναφοράς & Τι διατηρεί & Κύριος μηχανισμός ενημέρωσης \\
\hline
\ENG{Redis} \texttt{\ENG{LaunchRuntimeState}} & Προσωρινό έγκυρο αποθετήριο εκτέλεσης & Ενεργή κατάσταση συνεδρίας, τελευταία ακολουθία, κανονικοποιημένη πρόοδος, ακατέργαστη κατάσταση & \texttt{\ENG{CreateLaunchCommandHandler}}, \texttt{\ENG{RuntimeOrchestrator}} \\
\hline
\ENG{Postgres} \texttt{\ENG{attempts}} & Κύριο αποθετήριο προσπαθειών & Τελική πρόοδος, βαθμολογία, χρόνος, θέση, κατάσταση προσπάθειας & \texttt{\ENG{RuntimeFlushService}}, υπηρεσίες καταχώρισης \\
\hline
\ENG{Postgres} \texttt{\ENG{attempt\_runtime\_snapshots}} & Κύριο ιστορικό κατάστασης εκτέλεσης & Ακατέργαστο \ENG{runtime JSON}, κανονικοποιημένη πρόοδος, λόγος καταγραφής & \texttt{\ENG{RuntimeFlushService}} \\
\hline
\ENG{MinIO} \ENG{ZIP objects} & Κύριο αποθετήριο πακέτων περιεχομένου & Αρχικό συμπιεσμένο πακέτο του μαθήματος & \texttt{\ENG{ImportCourseCommandHandler}} \\
\hline
\ENG{Player cache} & Υποστηρικτική προσωρινή μνήμη, όχι κύριο αποθετήριο & Αποσυμπιεσμένα αρχεία για σερβίρισμα \ENG{SCO} & \texttt{\ENG{ContentCacheService.ensureExtracted()}} \\
\hline
\ENG{Browser in-memory state} & Εφήμερη κατάσταση πελάτη & Μεταβολές \ENG{SetValue} πριν από το επόμενο \ENG{commit} & \ENG{JavaScript} της σελίδας φιλοξενίας και \ENG{browser adapter} \\
\hline
\end{tabular}
\end{adjustbox}
\endgroup

\end{table}
